\section{Data, enclosing cuts, and the metric mass}\label{sec:setup}

We use the cosmological constant $\Lambda=-3$ and units $G=c=1$.
All initial data and boundaries are smooth. Norms at infinity are tensor
norms in the hyperbolic metric, rather than norms of coordinate components.

\begin{definition}[Initial-data exterior]\label{def:setup-data}
An initial-data exterior is a connected orientable three-manifold
$\Omega$ with nonempty compact smooth boundary $S$, a Riemannian metric
$g$, and a symmetric covariant tensor $K$. The metric space
$(\Omega,g)$, with its boundary included, is complete and has exactly one end.
Outside a compact set there is a chart $(r_0,\infty)\times\mathbb S^2$ with
\begin{equation}\label{eq:setup-background}
 b=\frac{dr^2}{1+r^2}+r^2\sigma,
 \qquad \eta=\operatorname{arsinh}r,
\end{equation}
where $\sigma$ is the unit round metric. For some $0<\alpha<1$ and
$3/2<\tau<3$, assume
\begin{equation}\label{eq:setup-decay}
 g-b\in C^{2,\alpha}_\tau,
 \qquad K\in C^{1,\alpha}_\tau.
\end{equation}
The space $C^{k,\alpha}_\tau$ controls $e^{\tau\eta}$ times the
$b$-norm of every covariant derivative through order $k$, and the
corresponding weighted H\"older seminorms on $b$-unit balls.
The noncosmological constraint densities are
\begin{equation}\label{eq:setup-constraints}
 16\pi\mu=R_g+6+(\tr_gK)^2-\lvert K\rvert_g^2,
 \qquad 8\pi J=\divg_g\bigl(K-(\tr_gK)g\bigr).
\end{equation}
We require
\begin{equation}\label{eq:setup-dec}
 \mu\ge\lvert J\rvert_g,
 \qquad
 \int_\Omega\sqrt{1+r^2}
       \bigl(\mu+\lvert J\rvert_g+\lvert K\rvert_g^2\bigr)
       \dd V_g<\infty,
\end{equation}
where the positive weight is extended smoothly over the compact part.
No restriction is placed on the compact topology, the number of boundary
components, or their genera.
\end{definition}

If $\nu$ is the normal toward the designated exterior and infinity, set
\begin{equation}\label{eq:setup-expansion}
 H=\divg_{S}\nu,
 \qquad \theta_+=H+\tr_{TS}K.
\end{equation}
The boundary hypothesis is $\theta_+\le0$ on every component of $S$,
with this same future sign. No condition on the other null expansion is
imposed. A spacetime realization uses
$K(Y,Z)=\overline g(\overline\nabla_Y n,Z)$ for the future unit normal $n$.
A surface with $\theta_+=0$ is a marginally outer trapped surface (MOTS).

\begin{definition}[Full enclosing cuts]\label{def:setup-cut}
An admissible exterior region is a connected smooth codimension-zero
submanifold with boundary $D\subset\Omega$, closed in $\Omega$, whose
interior lies in $\operatorname{int}\Omega$ and which contains the entire
sufficiently distant end. An admissible cut is its full compact intrinsic
boundary $\Gamma=\partial D$, assumed smoothly embedded and two-sided.
Coincidence with $S$, or with some of its components, is allowed. Define
\begin{equation}\label{eq:setup-area}
 A_{\min}(S;g)=\inf_{\Gamma}\Area_g(\Gamma),
 \qquad r_A=\sqrt{\frac{A_{\min}(S;g)}{4\pi}}.
\end{equation}
We omit $g$ when it is fixed. The competitor $D=\Omega$ has intrinsic
boundary $S$; it has area $\Area_g(S)$, not zero. Admissible cuts need not be
minimal or trapped, and may be disconnected.
\end{definition}

There is no extension across $S$ in this definition. An enclosing
minimization or trapped-region construction used below is always carried
out in the given exterior with $S$ as an obstacle or an actual boundary.
In particular, changing the exterior region requires comparison of the
entire cut, including any retained components of $S$.

\begin{definition}[Hyperbolic metric four-flux]\label{def:setup-mass}
Let $x_1,x_2,x_3$ be the coordinate functions on $\mathbb S^2$, and put
\begin{equation}\label{eq:setup-potentials}
 V_0=\sqrt{1+r^2},\qquad V_i=rx_i\quad(1\le i\le3).
\end{equation}
For $e=g-b$ and a static potential $V$, define the one-form
\begin{equation}\label{eq:setup-flux-form}
 \mathbb U_b(V,e)
 =V\bigl(\divg_b e-d\tr_be\bigr)
    +(\tr_be)dV-e(\nabla_bV,\cdot).
\end{equation}
The four mass components are
\begin{equation}\label{eq:setup-flux}
 p_a(g)=\frac1{16\pi}\lim_{R\to\infty}
   \int_{\{r=R\}}\mathbb U_b(V_a,e)(\nu_b)\dd A_b.
\end{equation}
The target class requires these limits to be finite and future timelike:
$p_0>\sqrt{p_1^2+p_2^2+p_3^2}$. Its mass is
\begin{equation}\label{eq:setup-invariant-mass}
 m_{\AH}(g)=\sqrt{p_0^2-p_1^2-p_2^2-p_3^2}.
\end{equation}
For auxiliary component estimates we do not impose timelikeness.
\end{definition}

The entries $p_i$ are spatial components of the metric mass covector.
They are not ADM momentum or angular momentum. This paper makes no
invariance assertion under arbitrary changes of asymptotic spacetime
slice. The asymptotic data are $(b,0)$ with $\Lambda=-3$.

\begin{lemma}[Flux convergence and spatial covariance]\label{lem:setup-mass-covariance}
The decay and integrability in Definition~\ref{def:setup-data} imply
convergence of the four integrals in Equation~\ref{eq:setup-flux}.
The same limits are obtained on any smooth exhaustion whose boundary
escapes every compact set. Under a fixed hyperbolic background isometry
the covector $p$ transforms by the corresponding Lorentz transformation.
It is invariant, with this transformation understood, under admissible
asymptotic spatial chart changes: both charts must satisfy the
tensor-decay and weighted-integrability conditions of
Definition~\ref{def:setup-data}. In particular, a future-timelike $p$
can be balanced by a spatial change of hyperbolic chart so that
$(p_0,p_1,p_2,p_3)=(m_{\AH},0,0,0)$.
\end{lemma}

\begin{proof}
Each $V_a$ satisfies $\Hess_bV_a=V_ab$ and $\Delta_bV_a=3V_a$.
The scalar curvature linearization at $b$ is
\[
 (DR)_b(e)=\divg_b\divg_be-\Delta_b\tr_be+2\tr_be.
\]
Differentiating Equation~\ref{eq:setup-flux-form} gives the exact identity
\begin{equation}\label{eq:setup-scalar-divergence}
 \divg_b\mathbb U_b(V_a,e)=V_a(DR)_b(e).
\end{equation}
In a sufficiently distant end write
$R_g+6=(DR)_b(e)+Q(e)$, where
\[
 \lvert Q(e)\rvert
 \le C\bigl(\lvert e\rvert\lvert\nabla_b^2e\rvert
                +\lvert\nabla_be\rvert^2+\lvert e\rvert^2\bigr).
\]
The right side is $O(e^{-2\tau\eta})$. Since
$|V_a|\le C e^\eta$ and $dV_b=O(e^{2\eta})d\eta dA_\sigma$,
the product $V_aQ$ is integrable precisely with the sufficient condition
$2\tau>3$. Moreover,
\[
 |R_g+6|\le16\pi\mu+4|K|_g^2,
\]
so $V_a(R_g+6)$ is integrable by Equation~\ref{eq:setup-dec} and the
uniform equivalence of $g$ and $b$ in the end. Equation~\ref{eq:setup-scalar-divergence}
and the divergence theorem now make the fluxes Cauchy. The difference
between their values on two escaping homologous cut boundaries is
bounded by the integral of this absolutely integrable divergence over
the intervening tail. This also proves the asserted exhaustion
independence for boundaries enclosing the same compact region.

In the hyperboloid model the four functions in
Equation~\ref{eq:setup-potentials} are its ambient coordinate functions.
A hyperbolic isometry therefore acts linearly on their span by a Lorentz
matrix. The flux expression is linear in $V$ and natural under a
background isometry. Its transformed coordinate spheres are an escaping
exhaustion, to which the preceding argument applies. This proves the
transformation law directly. The decay and weighted integrability remain
valid because a fixed hyperbolic isometry changes $\eta$ by a bounded
amount. The more general admissible spatial chart covariance is the
metric mass invariance Theorem of Chru\'sciel--Herzlich
\cite[Proposition~2.2 and Theorem~2.3]{chruscielherzlich2003}: its quadratic
decay threshold and weighted scalar-integrability hypotheses have just
been checked. Only the background-isometry case is needed to balance
the mass in the numerical proof. Finally, elementary Lorentz linear
algebra sends every future-timelike covector to its rest covector.
\end{proof}

\begin{remark}[Conformal normalization]\label{rem:setup-conformal-flux}
In three dimensions the linear flux for a pure conformal perturbation is
\begin{equation}\label{eq:setup-conformal-flux}
 \mathbb U_b(V,4t b)=-8(V\,dt-t\,dV).
\end{equation}
Indeed, $\tr_b(4tb)=12t$, $\divg_b(4tb)=4dt$, and the last two terms
of Equation~\ref{eq:setup-flux-form} sum to $8t\,dV$.
For example, if $t=O_1(r^{-s})$, $g-b=O_1(r^{-\tau})$,
$2s>3$, and $s+\tau>3$, then
\[
 e^{4t}g-g-4tb=O_1(r^{-2s})+O_1(r^{-s-\tau}).
\]
Its flux against each static potential is
$O(r^{3-2s})+O(r^{3-s-\tau})$, which tends to zero.
Thus Equation~\eqref{eq:setup-conformal-flux} gives the exact limiting
mass change whenever those limits exist; weighted scalar-curvature
integrability supplies their existence as above.
\end{remark}

\begin{definition}[Connected horizon subclass]\label{def:setup-horizon}
The horizon subclass consists of the preceding data with $S$ connected,
$\theta_+(S)=0$, and the following two properties. First, no compact
smooth embedded two-sided enclosing surface contained entirely in
$\operatorname{int}\Omega$, possibly disconnected and oriented on each
component toward infinity, has $\theta_+\le0$ everywhere. This is
outermostness. Second,
$\Area_g(S)\le\Area_g(\Gamma)$ for every full cut $\Gamma$.
This is outer area minimization, and gives
$A_{\min}(S)=\Area_g(S)$.
\end{definition}

\begin{lemma}[Schwarzschild--AdS normalization]\label{lem:setup-model-mass}
For $M>0$, let $F_M(r)=1+r^2-2M/r$ and let $r_h$ be its positive root.
The metric $g_M=F_M^{-1}dr^2+r^2\sigma$ on $[r_h,\infty)$, with
its smooth minimal-boundary interpretation, has
$(p_0,p_1,p_2,p_3)=(M,0,0,0)$. A full smooth spacelike cross-section
of the future horizon in the regular extension of
\begin{equation}\label{eq:setup-model-spacetime}
 \overline g_M=-F_M\,dt^2+F_M^{-1}dr^2+r^2\sigma
\end{equation}
has area $4\pi r_h^2$, and
$M=(r_h+r_h^3)/2$.
\end{lemma}

\begin{proof}
In the radial $b$-orthonormal frame, $g_M-b$ has just one entry,
\[
 q=\frac{2M}{rF_M},
\]
in the radial direction. Thus
$\mathbb U_b(V_0,g_M-b)(\nu_b)=2V_0\coth\eta\,q$.
The normalized integral is $M(1+r^2)/F_M$, which tends to $M$.
The three other integrals vanish by the zero spherical averages of
$x_i$. On the horizon the degenerate induced metric is the pullback of
$r_h^2\sigma$ from the space of null generators. A full spacelike
cross-section projects diffeomorphically onto that sphere and therefore
has its area. The relation between $M$ and $r_h$ follows from $F_M(r_h)=0$.
\end{proof}

Throughout the proof, estimates denoted by $O_j$ include derivatives
through order $j$ and the local H\"older control used in the relevant
elliptic estimate. Hyperbolic unit-ball estimates are distinguished from
asymptotic symbol estimates with derivatives $r\partial_r$ and
$\partial_\omega$. Preliminary approximation parameters are fixed before
subsequent deformation parameters. The final numerical limits are
limits of inequalities; no regular limiting deformation metric is
assumed in the equality argument.
